% 自動生成（tools/mknext.py p2）。各章の道しるべ「この章の先に」から、後ろの巻の章を集めたもの（第13回）
\begin{longtblr}{colspec={Q[l,m,wd=58mm] X[l,m]}, width=\linewidth, rowhead=1, row{1}={bg=bkVol!14, font=\sffamily\bfseries\footnotesize},
  hlines={0.4pt,bkLine}, rowsep=1.5pt, cells={font=\footnotesize}}
この巻の章 & この道具を使う章 \\
{\bknextpill{熱-3}\ ふくらむ・縮む気体：理想気体の状態変化} & {\bknextpill{波-9}\ 音の大きさをはかる：音速・デシベル・騒音} \\
{\bknextpill{熱-4}\ 流れる気体の熱：エンタルピーとノズル} & {\bknextpill{波-9}\ 音の大きさをはかる：音速・デシベル・騒音} \\
{\bknextpill{熱-7}\ エンジンの中のサイクル：ガスサイクル} & {\bknextpill{電-12}\ モーターと変圧器のしくみ：電気機器} \\
{\bknextpill{熱-9}\ 冷やす・湿り気をはかる：冷凍と湿り空気} & {\bknextpill{波-4}\ ねらいどおりに動かす：フィードバック制御} \\
{\bknextpill{熱-10}\ わく温度を圧力で読む：自由エネルギーと相平衡} & {\bknextpill{電-4}\ 電気の流れと熱：電流と抵抗} \\
{\bknextpill{熱-11}\ 壁を通りぬける熱：熱伝導と熱抵抗} & {\bknextpill{波-4}\ ねらいどおりに動かす：フィードバック制御} \\
{\bknextpill{熱-13}\ 光で伝わる熱：ふく射と熱交換器} & {\bknextpill{波-12}\ 光の進み方：反射・屈折・レンズ\\ \bknextpill{波-13}\ 光の波：干渉・回折・偏光・散乱\\ \bknextpill{原-2}\ 光はつぶでもある：光電効果と光子} \\
{\bknextpill{熱-14}\ 温度の正体：気体分子運動論} & {\bknextpill{波-9}\ 音の大きさをはかる：音速・デシベル・騒音\\ \bknextpill{原-6}\ 壁をすりぬける：トンネル効果と調和振動子} \\
{\bknextpill{熱-15}\ 並び方を数える：統計力学の入口} & {\bknextpill{原-4}\ 原子が出す光：エネルギー準位とスペクトル\\ \bknextpill{原-6}\ 壁をすりぬける：トンネル効果と調和振動子\\ \bknextpill{原-8}\ 固体の中の電子：エネルギーバンド\\ \bknextpill{原-9}\ 一方通行の電気：pn接合とLED・太陽電池} \\
{\bknextpill{熱-16}\ エネルギーをどう選ぶか} & {\bknextpill{電-16}\ 電気を届ける・電気を変える\\ \bknextpill{原-15}\ 原子炉の中の中性子：核分裂と原子炉\\ \bknextpill{原-16}\ 原子の力とつきあう} \\
\end{longtblr}
